In this chapter, we introduce a fundamental idea of discrete mathematics and computer science: recursion. Essentially, recursion is a problem-solving strategy in which we replace the problem we are trying to solve with a smaller, hopefully easier, version of the same problem.  We continue this replacement with smaller and smaller versions of the problem until we eventually reach a version of the problem that is small enough to know the answer.  

We start our exploration of recursions by looking at recursively-defined sequences in \Cref{sec_recursively_defined_sequences} {\secrecseq}.  Next, in \Cref{sec_math_ind} {\secmathind}, we introduce proof by mathematical induction, a proof format that is particularly useful for verifying that recursively-defined structures have explicit properties.  In \Cref{sec_matrix} {\secmatrix}, we take a brief excursion to visit matrices and finite groups, which also gives us the opportunity to practice proofs by mathematical induction.  \Cref{sec_fibonacci_higher} {\secfibhigher} we explore the Fibonacci sequence and other examples of higher order recursions and introduce an extension of proof by mathematical induction.  Last, we investigate other recursively-defined structures, including sets and graphs, and how to count with recursions in \Cref{sec_rec_def_structures} {\secrecdefn}.


